\subsection{Variational mechanics and the conservative--forced split}
\label{mech:variational}
The residual-compensation construction uses one mechanical object
throughout reduction, identification, discretization and actuation.
We first state the general variational entrance, then specialize it to
a two-mode magnetic beam with known kinetic energy. The coupled
magnetic potential remains in the control target; the learned terms
represent a conservative potential residual and an effective damping
rate. The parameter choices below specify this object and its physical
input scales.

The mechanical starting point is a rod with centreline $r(s,t)$ and
director frame $D(s,t)\in SO(3)$, where $s\in[0,L]$ is reference
arclength. A configuration map $(r(s,q),D(s,q))$ satisfying the essential
boundary conditions and any imposed rod constraints determines both
translational and rotational kinetic energies \citep{simo,antman}.
After spatial reduction, write
\begin{equation}
 T=\tfrac12\dot q^{\mathsf T}M(q)\dot q,\qquad V=V(q),
 \label{mech:general-energy}
\end{equation}
with $M(q)=M(q)^{\mathsf T}\succ0$. Elasticity and a prescribed static
magnetic field contribute to $V$; modulation of that field supplies work.
For a permanent dipole, the interaction potential $-m\cdot B$ gives
$\tau=m\times B$ and $F=\nabla(m\cdot B)$, with orientation held fixed
in the spatial derivative \citep{abbott}.

Let $c_{\rm geo}(q,\dot q)$ be the inertial term generated by derivatives
of $M$, let $\mathcal D_N(q,\dot q)$ be a Rayleigh dissipation potential,
and let $B_u(q)u$ denote the generalized control force. The variational
equations and their Legendre transform are
\begin{align}
 M\ddot q+c_{\rm geo}+\nabla V
 &=Q_{\rm drv}-\partial_{\dot q}\mathcal D_N+B_u u,
 \label{mech:general-lagrange}\\
 p&=M(q)\dot q,\qquad
 H=\tfrac12p^{\mathsf T}M(q)^{-1}p+V,\nonumber\\
 \dot q&=H_p,\qquad v=H_p,\nonumber\\
 \dot p&=-H_q-\partial_v\mathcal D_N(q,v)+Q_{\rm drv}+B_u u.
 \label{mech:general-hamiltonian}
\end{align}
In particular, $(H_q)_i$ contains
$p^{\mathsf T}(\partial_{q_i}M^{-1})p/2$.
The geometric inertia is not an additional dissipative force.
Fluid forcing, moving-boundary work and contact losses can be retained
in $Q_{\rm drv}$ when their models are specified. The derivation in
Supplement~\ref{sm:general} retains these distinctions, including the
fact that a moving fluid can supply mechanical energy. Analytic
persistence results below are restricted to a smooth, contact-free
neighbourhood.

\subsection{A literature-informed physical benchmark}
\label{mech:benchmark}
To obtain a reproducible instance, we specialize this formulation to a
homogeneous circular cantilever in a plane, with a rigid magnetic core
at its distal end. The model has a fixed root and no transverse gravity,
external flow, contact or fluid added mass. This is a gravity-free
mechanical benchmark motivated by remote manipulation, including
microgravity settings; it is not a patient-specific catheter model.
Euler--Bernoulli kinematics neglect shear deformation and distributed
rotary inertia. The tip rotation is approximated by
$\vartheta=w_s(L,t)$; retaining its trigonometric magnetic potential
does not remove this small-slope approximation.

We denote modal damping ratios by $\zeta_j^{\mathrm{d}}$.
Table~\ref{mech:parameters} separates measured material reference
values, nominal magnetic specifications and design choices. The Sylgard
184 reference uses approximately $10$ wt\% curing agent, cured at
$80\,^{\circ}\mathrm C$ for $2$ h \citep{cho2021pdms}. The reported
$E$ and density are approximate reference values for that preparation,
not measurements of the present cylindrical specimen. A $10$ wt\%
curing-agent fraction corresponds to approximately $9{:}1$ base to
curing agent by mass, rather than $10{:}1$.

\begin{table*}[t]
\centering\small
\caption{Inputs to the revised mechanical benchmark. Values designated as
design choices or scenarios are not experimental estimates.}
\label{mech:parameters}
\setlength{\tabcolsep}{5pt}
\begin{tabular}{p{0.27\textwidth}p{0.22\textwidth}p{0.44\textwidth}}
\toprule
Quantity & Value & Basis and interpretation\\
\midrule
Beam length $L$, radius $r$ & $80$ mm, $0.6$ mm & Declared geometry of the uniform solid beam.\\
Young's modulus $E$ & approximately $2.1$ MPa & Sylgard 184 material reference, with the preparation stated in the text \citep{cho2021pdms}.\\
Mass density $\rho$ & approximately $1028\,\mathrm{kg\,m^{-3}}$ & Same material reference \citep{cho2021pdms}; no same-batch specimen calibration is asserted.\\
Bare magnetic-core edge $a_m$ & $0.250$ mm & Declared cubic-core design; coating and adhesive are excluded.\\
Remanent flux density $B_r$ & nominal $1.425$ T & Midpoint of the N50 grade range $1.40$--$1.45$ T \citep{smN50}.\\
Core density $\rho_m$ & nominal $7500\,\mathrm{kg\,m^{-3}}$ & Selected at the manufacturer's stated lower specification; not an individual-core density measurement \citep{smN50}.\\
Axial bias $B_0$ & $30$ mT & Design choice; published OctoMag experiments provide a $40$ mT vector-field capability comparison \citep{lussi2021cvs}.\\
First modal damping ratio $\zeta_1^{\mathrm{d}}$ & $0.025$ & Effective scenario, defined relative to the biased frequency including tip inertia.\\
Effective centreline clearance $c$ & $4$ mm & Design comparison, after subtracting the rod radius from the corridor half-width.\\
\bottomrule
\end{tabular}
\end{table*}

For a uniformly magnetized hard-magnetic core we use the engineering
approximation $\mu\simeq B_ra_m^3/\mu_0$, with
$\mu_0\simeq4\pi10^{-7}\,\mathrm{H\,m^{-1}}$ \citep{kjMoment}.
The nominal dipole moment and inertial properties are
\begin{align}
 \mu&=1.77184214\times10^{-5}\,\mathrm{A\,m^2},\nonumber\\
 m_t=\rho_ma_m^3&=1.171875\times10^{-7}\,\mathrm{kg},\nonumber\\
 J_t=m_ta_m^2/6&=1.220703125\times10^{-15}\,\mathrm{kg\,m^2}.
 \label{mech:core-properties}
\end{align}
The core's centre of mass is idealized at the beam endpoint. An
eccentric installation would introduce inertial coupling and altered
load moment arms. The magnetic-moment estimate does not identify the
magnetization history, coating or manufacturing tolerances of a device.
These assumptions define a consistent numerical object without
combining the material properties of different composite robots.

\subsection{Boundary inertia, finite elements and stored energy}
\label{mech:finite-elements}
Set $A_c=\pi r^2$, $I=\pi r^4/4$, and use the effective damping model
$C=\alpha M_f$. For the associated continuous beam, the interior
equation and clamped boundary are
\begin{align}
 \rho A_cw_{tt}+\alpha\rho A_cw_t+EIw_{ssss}&=0,\nonumber\\
 w(0,t)=w_s(0,t)&=0.
 \label{mech:beam}
\end{align}
Writing $F$ for a transverse tip force and
$\tau_B=\mu B_y\cos\vartheta$ for the active tip torque, the natural
endpoint balances become
\begin{align}
 -EIw_{sss}(L)&=F-m_t[w_{tt}(L)+\alpha w_t(L)],\nonumber\\
 EIw_{ss}(L)&=\tau_B-\mu B_0\sin\vartheta\nonumber\\
 &\quad-J_t(\ddot\vartheta+\alpha\dot\vartheta).
 \label{mech:endpoint-balance}
\end{align}
Thus the terminal rotation remains free. The endpoint damping terms are
part of the declared $C=\alpha M_f$ approximation, not independently
measured losses in the magnetic core.

Let $q$ collect the unconstrained Hermite displacement and slope
coordinates, and let $e_w,e_\vartheta$ select their tip values. The weak
form gives
\begin{align}
 M_f&=M_{\rm beam}+m_te_we_w^{\mathsf T}
                         +J_te_\vartheta e_\vartheta^{\mathsf T},\nonumber\\
 M_f\ddot q+\alpha M_f\dot q+K_{\rm beam}q
 &=e_wF+e_\vartheta(\tau_B-\mu B_0\sin\vartheta),\nonumber\\
 \vartheta&=e_\vartheta^{\mathsf T}q.
 \label{mech:fe-plant}
\end{align}
With $p=M_f\dot q$ and $E_0=\mu B_0$, the nominal physical Hamiltonian is
\begin{align}
 H_0(q,p)&=\tfrac12p^{\mathsf T}M_f^{-1}p
             +\tfrac12q^{\mathsf T}K_{\rm beam}q\nonumber\\
 &\quad+E_0[1-\cos\vartheta].
 \label{mech:fe-energy}
\end{align}
Here $E_0=5.31552642\times10^{-7}\,\mathrm J$.
The bias stores energy, while $F$ and $B_y$ supply work through distinct
physical ports. Their power balance is
\begin{equation}
 \dot H_0=-\alpha\dot q^{\mathsf T}M_f\dot q
              +F\dot w(L)+\tau_B\dot\vartheta.
 \label{mech:physical-power}
\end{equation}

The computation uses $32$ consistent-mass cubic-Hermite elements and
$64$ free nodal coordinates. The independent endpoint characteristic
equation and $16/32/64$-element comparison in
Supplement~\ref{sm:fe-check} verify the first two frequencies to relative
errors $9.02\times10^{-9}$ and $3.29\times10^{-7}$ at $32$ elements.
This checks linear discretization accuracy; the subsequent two-mode
restriction remains a separate model reduction.

\subsection{Mass normalization and canonical scaling}
\label{mech:canonical}
Linearizing the static magnetic term gives
\begin{align}
 K_b&=K_{\rm beam}+E_0e_\vartheta e_\vartheta^{\mathsf T},\nonumber\\
 K_b\Phi_n&=M_f\Phi_n\Omega_n^2,\qquad
 \Phi_n^{\mathsf T}M_f\Phi_n=I_n,
 \label{mech:modal-bias}
\end{align}
where $\Omega_n=\diag(\omega_1,\ldots,\omega_n)$.
With $q=\Phi_na$, the modal momentum is $\pi=\dot a$.
The modal mass is $I_n$ only after this normalization.
Set $\omega_0=\omega_1$, $\tau=\omega_0t$,
$\nu_j=\omega_j/\omega_0$ and $\mathsf N=\diag(\nu_j)$.
The balanced canonical scaling is
\begin{align}
 a_j&=S_jX_j,\qquad S_j=\sqrt{\frac{E_0}{\omega_0\omega_j}},\nonumber\\
 \pi_j&=\sqrt{\frac{E_0\omega_j}{\omega_0}}Y_j,
 \qquad D_s=\diag(S_j).
 \label{mech:scaling}
\end{align}
The position--momentum scale product is $E_0/\omega_0$.
Together with energy $H_0/E_0$ and time $\tau$, this gives canonical
dimensionless equations, rather than a unit-mass replacement of $M_f$.
For the Hermite interpolation row $N_w(s)$, define
\begin{equation}
 \ell(s)^{\mathsf T}=N_w(s)\Phi_nD_s,\qquad
 d=D_s\Phi_n^{\mathsf T}e_\vartheta.
 \label{mech:shape-maps}
\end{equation}
Then $w(s)=\ell(s)^{\mathsf T}X$, $\vartheta=d^{\mathsf T}X$, and
\begin{align}
 h_0(X,Y)&=h_2(X,Y)+g(d^{\mathsf T}X),\nonumber\\
 h_2&=\tfrac12(X^{\mathsf T}\mathsf NX+Y^{\mathsf T}\mathsf NY),\nonumber\\
 g(s)&=1-\cos s-\tfrac12s^2.
 \label{mech:nominal-energy}
\end{align}
The quadratic subtraction avoids counting the bias stiffness twice.
For the two-mode instance,
\begin{align}
 \omega_0&=7.84000916\,\mathrm{rad\,s^{-1}},\nonumber\\
 \nu&=(1,6.04585730)^{\mathsf T},\nonumber\\
 d&=(0.30722439,0.45042632)^{\mathsf T},\nonumber\\
 \ell(L)&=(0.0189673443,0.00775347614)^{\mathsf T}\,\mathrm m.
 \label{mech:reduced-values}
\end{align}
The exported coefficients retain more digits than displayed here.
Additional digits identify the numerical benchmark and do not imply
corresponding measurement precision.
The effective mass in $(X,Y,\tau)$ is $\mathsf N^{-1}$.
Table~\ref{mech:dimension-table} collects the coordinate levels used
by the residual controller and its discretization. The $32$-element
model has $64$ configuration coordinates and a $128$-dimensional phase
space. Restricting its mass-normalized modal family to $n=2$ gives
$I_2$, while the subsequent balanced scaling gives
$M_{\rm eff}=\mathsf N^{-1}$. Each mass belongs to its stated
coordinates. For a general $n$-mode model, the phase-flow Jacobian
$D\Phi$ has size $2n\times2n$, and the joint endpoint transformation
has size $4n\times4n$; in the controlled two-mode system these sizes
are $4\times4$ and $8\times8$. Supplement~\ref{ss:projection}
gives the endpoint projection used by the generating function.

\begin{table*}[t]
\centering\small
\caption{Coordinate levels of the same mechanical benchmark.
Mass normalization and balanced scaling are separate transformations.}
\label{mech:dimension-table}
\setlength{\tabcolsep}{5pt}
\begin{tabular}{p{.22\textwidth}p{.24\textwidth}p{.22\textwidth}p{.23\textwidth}}
\toprule
Level & Coordinates and time & Kinetic mass & Phase-state and tangent sizes\\
\midrule
Complete $32$-element beam & $q,p\in\mathbb R^{64}$, physical $t$ & $M_f\in\mathbb R^{64\times64}$, including $m_t,J_t$ & $128$; $D\Phi\in\mathbb R^{128\times128}$\\
Mass-normalized $n$ modes & $a,\pi\in\mathbb R^n$, physical $t$ & $I_n$; $I_2$ when $n=2$ & $2n$; $D\Phi\in\mathbb R^{2n\times2n}$\\
Balanced two-mode controller & $X,Y\in\mathbb R^2$, $\tau=\omega_0t$ & $M_{\rm eff}=\mathsf N^{-1}$ & $4$; $D\Psi_\Delta\in\mathbb R^{4\times4}$\\
\bottomrule
\end{tabular}
\end{table*}

The damping scenario is
$\alpha=2\zeta_1^{\mathrm{d}}\omega_0=0.3920004581\,\mathrm{s^{-1}}$ and
$\gamma=\alpha/\omega_0=0.05$; it does not assign the same damping ratio
to every mode. Other damping ratios are explored as declared scenarios,
without substituting material loss factors measured under different
conditions for modal damping.

\subsection{Complex coordinates, dissipation and physical inputs}
\label{mech:complex}
Let $x=(X,Y)$ and $J=\left(\begin{smallmatrix}0&I\\-I&0\end{smallmatrix}\right)$.
The invertible complex chart is
\begin{equation}
 \zeta=(X+\mathrm iY)/\sqrt2,\qquad
 \chi=(X-\mathrm iY)/\sqrt2.
 \label{mech:complex-chart}
\end{equation}
On real phase space $\chi=\bar\zeta$. Treating $\zeta,\chi$ as
independent complex variables gives the entire continuation
\begin{equation}
 h_0^{\mathbb C}(\zeta,\chi)
 =\sum_j\nu_j\zeta_j\chi_j
       +g\!\left(d^{\mathsf T}(\zeta+\chi)/\sqrt2\right).
 \label{mech:complex-energy}
\end{equation}
With the convention
$\{\zeta_j,\bar\zeta_k\}=-\mathrm i\delta_{jk}$,
\begin{align}
 \partial_{\bar\zeta}&=(\partial_X+\mathrm i\partial_Y)/\sqrt2,\nonumber\\
 \zeta'&=-\mathrm i\partial_\chi h^{\mathbb C},\qquad
 \chi'=\mathrm i\partial_\zeta h^{\mathbb C}.
 \label{mech:wirtinger}
\end{align}
Analytic continuation in two independent complex arguments is the object
used in analytic KAM theory; a real energy need not be holomorphic in
$\zeta$ alone.

A complex Hamiltonian neural representation supplies a real energy
readout $h_\eta$ and a dimensionless dissipation readout
$\mathcal D_\eta(X,v)$, obtained by dividing physical Rayleigh
dissipation by $E_0\omega_0$. Its dynamics are defined explicitly by
\begin{align}
 X'&=v=\partial_Yh_\eta,\nonumber\\
 Y'&=-\partial_Xh_\eta-\partial_v\mathcal D_\eta(X,v)+U.
 \label{mech:decoder}
\end{align}
The velocity derivative is taken before substituting
$v=\partial_Yh_\eta$. For example,
$\mathcal D_\eta=\widehat\gamma v^{\mathsf T}\mathsf N^{-1}v/2$,
$\widehat\gamma\ge0$, gives
\begin{equation}
 h_\eta'=-v^{\mathsf T}\partial_v\mathcal D_\eta+v^{\mathsf T}U,
 \label{mech:power}
\end{equation}
with $\partial_\tau h_\eta$ added for time-dependent parameters.
Thus $h_\eta+\mathrm i\mathcal D_\eta$ is a pair of decoded readouts,
distinct from the analytic continuation $h_\eta^{\mathbb C}$.
In complex notation the nonconservative term is
$\mathrm i[-\partial_v\mathcal D_\eta+U]/\sqrt2$.

For the fixed-kinetic two-mode model, put
$b_j=e_w^{\mathsf T}\phi_j$, $t_j=e_\vartheta^{\mathsf T}\phi_j$.
The physical input conversion is
\begin{align}
 A_{\rm mag}&=\begin{pmatrix}b_1&t_1\\b_2&t_2\end{pmatrix},
 &\Lambda&=\omega_0^2D_s\mathsf N=E_0D_s^{-1},\nonumber\\
 \Lambda U&=A_{\rm mag}\binom{F}{\tau_B}.
 \label{mech:modal-ports}
\end{align}
A uniform transverse field alone gives
$U=(B_y/B_0)d\cos\vartheta$, a single input direction. Independent
field and gradient authority is needed to realize a general two-vector
$U$. Section~\ref{sec:application} applies this interface to the learned
residual-compensation law and follows the same law into higher-mode
plants. A stronger realization of general target Hamiltonians,
including learned momentum terms, is stated separately in
Supplement~\ref{sm:exact-lift}; the tracking extension uses the shared
interface in Supplementary Section~\ref{sec:supp-tracking}.
